%% MCM/ICM Paper Template — based on mcmthesis
%% Compile: xelatex main.tex (run 3 times for cross-references + BibTeX)
%%
%% mcmthesis package: https://github.com/latexstudio-org/mcmthesis
%% If not installed: tlmgr install mcmthesis
%%
%% !! FILL IN: tcn (Team Control Number) and problem (A/B/C/D/E/F) !!

\documentclass{mcmthesis}

%% ── Team Configuration ─────────────────────────────────────────────────
\mcmsetup{
    tcn          = {0000000},     %% <── 7-digit Team Control Number
    problem      = {A},           %% <── Problem choice: A B C D E F
    sheet        = true,
    titleinsheet = true,
    keywordsinsheet = true,
    titlepage    = false,
    abstract     = true,
}

%% ── Packages ────────────────────────────────────────────────────────────
\usepackage{amsmath, amssymb, amsthm}
\usepackage{graphicx}
\usepackage{subcaption}
\usepackage{booktabs}
\usepackage{longtable}
\usepackage{multirow}
\usepackage{array}
\usepackage{float}
\usepackage{hyperref}
\usepackage{xcolor}
\usepackage{listings}
\usepackage{enumerate}
\usepackage{enumitem}
\usepackage{caption}
\usepackage{setspace}
\usepackage{tikz}
\usepackage{pgfplots}
\usepackage[ruled,vlined,linesnumbered]{algorithm2e}
% natbib 要求 author-year 格式的 \citep/\citet，但本模板正文用标准 \cite +
% plain \bibitem，直接加载会在编译时报 "Bibliography not compatible with
% author-year citations" 致命错误——同 CUMCM 模板的修复，不加载 natbib。
\usepackage{appendix}
\usepackage{microtype}           % better text justification
\usepackage{siunitx}             % units formatting
\usepackage{cleveref}            % smart cross-references

\pgfplotsset{compat=1.18}
\usetikzlibrary{shapes.geometric, arrows.meta, positioning, calc, fit}
\hypersetup{colorlinks=true, linkcolor=blue!70!black, citecolor=green!50!black,
            urlcolor=blue!70!black}

%% ── Code Style ──────────────────────────────────────────────────────────
\definecolor{backcolour}{rgb}{0.97,0.97,0.97}
\definecolor{codeblue}{rgb}{0.13,0.29,0.53}
\definecolor{codegray}{rgb}{0.5,0.5,0.5}
\definecolor{codepurple}{rgb}{0.58,0,0.82}

\lstdefinestyle{pythonstyle}{
    backgroundcolor=\color{backcolour},
    commentstyle=\color{codegray}\itshape,
    keywordstyle=\color{codeblue}\bfseries,
    stringstyle=\color{codepurple},
    basicstyle=\ttfamily\footnotesize,
    breaklines=true, captionpos=b,
    numbers=left, numbersep=5pt,
    numberstyle=\tiny\color{codegray},
    showstringspaces=false, tabsize=4,
    frame=single, rulecolor=\color{black!20},
    language=Python
}

%% ── Custom Commands ─────────────────────────────────────────────────────
\newtheorem{assumption}{Assumption}
\newtheorem{definition}{Definition}
\newtheorem{theorem}{Theorem}
\newtheorem{lemma}{Lemma}

\newcommand{\vect}[1]{\boldsymbol{#1}}
\newcommand{\norm}[1]{\left\lVert#1\right\rVert}
\newcommand{\abs}[1]{\left|#1\right|}
\newcommand{\E}[1]{\mathbb{E}\!\left[#1\right]}
\newcommand{\Prob}[1]{\mathbb{P}\!\left(#1\right)}
\DeclareMathOperator*{\argmin}{arg\,min}
\DeclareMathOperator*{\argmax}{arg\,max}

\setstretch{1.15}

%% ── Paper Title & Keywords ───────────────────────────────────────────────
\title{[Paper Title — Concise and Descriptive]}

\begin{document}

%% ══════════════════════════════════════════════════════════════════
%%  SUMMARY  (most critical section — 1 page maximum)
%%  mcmthesis renders this on the Summary Sheet automatically.
%% ══════════════════════════════════════════════════════════════════
\begin{abstract}

% Paragraph 1: Problem Context (2–3 sentences)
[State the problem context. What is being optimized/predicted/analyzed?
What are the key challenges?]

% Paragraph 2: Models Developed (4–6 sentences)
To address these challenges, we develop [N] interconnected models.
For \textbf{Task 1}, we employ [Model Name], formulated as
$\min_{x} f(x)$ subject to $g(x) \leq 0$, solved via [algorithm].
For \textbf{Task 2}, we construct a [Model Name] to [objective].
[Continue for each task.]

% Paragraph 3: Key Results (bullet form preferred)
Our analysis yields the following key findings:
\begin{itemize}[nosep]
    \item \textbf{Task 1}: [Specific quantitative result, e.g., optimal value = X.XX with 94.2\% confidence]
    \item \textbf{Task 2}: [Specific result]
    \item \textbf{Sensitivity}: The model remains robust to ±20\% perturbations in [key parameter].
\end{itemize}

% Paragraph 4: Highlight sentence
[The most impressive result in one sentence — make judges want to read on.]

\end{abstract}

\keywords{[keyword1]; [keyword2]; [keyword3]; [keyword4]; [keyword5]}

\maketitle

%% ── Table of Contents ────────────────────────────────────────────────────
%% AutoMCM_SOP.md §17：COMAP 官方规则（Contest_AI_Policy / instructions.php）——
%% 全篇合计（Summary Sheet + 正文 + 参考文献 + 目录 + 附录代码）硬上限 25 页，
%% 目录本身计入这 25 页；tocdepth 限到 2（section+subsection，不列
%% subsubsection）把目录控制在一页左右。若用到 AI，末尾追加的
%% "Report on Use of AI Tools" 单独一节不计入 25 页（见文末对应位置）。
\setcounter{tocdepth}{2}
\tableofcontents
\newpage

%% ══════════════════════════════════════════════════════════════════
%%  1. INTRODUCTION
%% ══════════════════════════════════════════════════════════════════
\section{Introduction}

\subsection{Background and Motivation}

[Macro-level context: why does this problem matter in the real world?
Cite 2–3 sources here to establish credibility.]

\subsection{Literature Review}

[Summarize existing approaches to this class of problems in 1–2 paragraphs.
Identify the gap your approach addresses. Cite ≥5 papers.]

Recent work by \cite{ref1} demonstrated that [key insight].
However, [limitation of existing work], which motivates the present study.

\subsection{Problem Restatement}

Given [input data/parameters], we seek to:
\begin{enumerate}
    \item [Task 1: formal mathematical statement]
    \item [Task 2: formal mathematical statement]
    \item [Task 3: formal mathematical statement]
\end{enumerate}

\subsection{Overview of Our Approach}

Figure~\ref{fig:flowchart} illustrates the overall modeling framework.

\begin{figure}[H]
\centering
\begin{tikzpicture}[
    box/.style={rectangle, rounded corners=4pt, draw=blue!60, fill=blue!8,
                text width=3cm, align=center, minimum height=0.85cm, font=\small},
    arrow/.style={-Stealth, thick, blue!60}
]
    \node[box] (data)   {Data Acquisition \& EDA};
    \node[box, right=1.4cm of data]  (m1)  {Model 1\\ [Task 1]};
    \node[box, right=1.4cm of m1]    (m2)  {Model 2\\ [Task 2]};
    \node[box, right=1.4cm of m2]    (out) {Results \&\\ Sensitivity};

    \draw[arrow] (data) -- (m1);
    \draw[arrow] (m1)   -- (m2);
    \draw[arrow] (m2)   -- (out);
\end{tikzpicture}
\caption{Overall modeling framework.}
\label{fig:flowchart}
\end{figure}

%% ══════════════════════════════════════════════════════════════════
%%  2. ASSUMPTIONS AND JUSTIFICATIONS
%% ══════════════════════════════════════════════════════════════════
\section{Assumptions and Justifications}

To ensure mathematical tractability while preserving fidelity to the real-world scenario, we adopt the following assumptions:

\begin{assumption}
    [Assumption statement.]
    \textit{Justification:} [Physical/empirical reason. Cite if possible.]
\end{assumption}

\begin{assumption}
    [Assumption statement.]
    \textit{Justification:} [...]
\end{assumption}

% Add 5–8 assumptions total

%% ══════════════════════════════════════════════════════════════════
%%  3. NOTATION
%% ══════════════════════════════════════════════════════════════════
\section{Notation}

Table~\ref{tab:notation} summarizes the principal mathematical notation used throughout this paper.

\begin{table}[H]
\centering
\caption{Principal notation.}
\label{tab:notation}
\begin{tabular}{p{2.2cm} p{7cm} p{2.5cm}}
\toprule
\textbf{Symbol} & \textbf{Definition} & \textbf{Unit / Range} \\
\midrule
$x_i$       & [definition]          & [unit] \\
$y_j$       & [definition]          & [unit] \\
$\alpha$    & [parameter meaning]   & $[0, 1]$ \\
$T$         & Time horizon          & days \\
$n$         & Number of [entities]  & $\mathbb{Z}^+$ \\
% Add more rows as needed
\bottomrule
\end{tabular}
\end{table}

%% ══════════════════════════════════════════════════════════════════
%%  4. MODEL DEVELOPMENT
%% ══════════════════════════════════════════════════════════════════
\section{Model Development}

% ────────── 4.1 Sub-problem 1 ──────────────────────────────────
\subsection{Task 1: [Model Name] for [Objective]}

\subsubsection{Mechanism Analysis}

[Describe the physical/mathematical mechanism. Why does this model structure fit the problem?]

\subsubsection{Mathematical Formulation}

We formulate Task~1 as a [type of problem]:

\begin{equation}
    \min_{\vect{x} \in \mathcal{X}} \quad f(\vect{x}) = \sum_{i=1}^{n} w_i \phi(x_i)
    \label{eq:obj_task1}
\end{equation}
subject to:
\begin{align}
    g_i(\vect{x}) &\leq 0, \quad i = 1, \ldots, m   \label{eq:ineq} \\
    h_j(\vect{x}) &=    0, \quad j = 1, \ldots, p   \label{eq:eq}   \\
    \vect{x}      &\geq \vect{0}                     \label{eq:nn}
\end{align}

where $\phi(\cdot)$ denotes [function meaning] and $\mathcal{X}$ is [feasible set description].

\subsubsection{Algorithm Design}

\begin{algorithm}[H]
\SetAlgoLined
\KwIn{[Inputs]}
\KwOut{[Outputs]}
Initialize [parameters]\;
\While{not converged}{
    Compute [step]\;
    Update [variables]\;
    Check termination criterion\;
}
\Return{[result]}\;
\caption{[Algorithm Name] for Task 1}
\end{algorithm}

\subsubsection{Results and Analysis}

Table~\ref{tab:result_task1} and Figure~\ref{fig:result_task1} present the key results.

\begin{table}[H]
\centering
\caption{Task 1 results summary.}
\label{tab:result_task1}
\begin{tabular}{lccc}
\toprule
\textbf{Metric} & \textbf{Value} & \textbf{Unit} & \textbf{Benchmark} \\
\midrule
[Metric 1] & [value] & [unit] & [comparison] \\
[Metric 2] & [value] & [unit] & [comparison] \\
\bottomrule
\end{tabular}
\end{table}

\begin{figure}[H]
    \centering
    \includegraphics[width=0.75\textwidth]{images/fig01_task1_result.png}
    \caption{[Figure caption with specific observation about the result.]}
    \label{fig:result_task1}
\end{figure}

[Interpret the results quantitatively. What does Figure~\ref{fig:result_task1} show?
Compare to any benchmarks. Discuss implications.]

% ────────── 4.2 Sub-problem 2 ──────────────────────────────────
\subsection{Task 2: [Model Name] for [Objective]}

[Same structure: Mechanism → Formulation → Algorithm → Results]

% ────────── 4.3 Sub-problem 3 ──────────────────────────────────
\subsection{Task 3: [Model Name] for [Objective]}

[Same structure. Add/remove subsections as needed.]

%% ══════════════════════════════════════════════════════════════════
%%  5. SENSITIVITY ANALYSIS
%% ══════════════════════════════════════════════════════════════════
\section{Sensitivity Analysis}

To assess the robustness of our models, we conduct a systematic sensitivity analysis on the key parameters identified in \cref{tab:notation}.

\subsection{Parameter Perturbation Study}

We vary [parameter $\alpha$] by $\{-50\%, -20\%, -10\%, 0\%, +10\%, +20\%, +50\%\}$ while holding other parameters at their baseline values.

\begin{table}[H]
\centering
\caption{Sensitivity of [objective] to perturbations in $\alpha$.}
\label{tab:sensitivity}
\begin{tabular}{cccc}
\toprule
$\Delta\alpha$ & Objective Value & Change (\%) & Feasible? \\
\midrule
$-50\%$ & & & \\
$-20\%$ & & & \\
$-10\%$ & & & \\
$0\%$   & & 0.0 & \checkmark \\
$+10\%$ & & & \\
$+20\%$ & & & \\
$+50\%$ & & & \\
\bottomrule
\end{tabular}
\end{table}

[Interpret: the model is [robust/sensitive] to changes in $\alpha$, with a [X\%] change in objective for a [Y\%] parameter perturbation.]

%% ══════════════════════════════════════════════════════════════════
%%  6. STRENGTHS AND WEAKNESSES
%% ══════════════════════════════════════════════════════════════════
\section{Strengths and Weaknesses}

\subsection{Strengths}
\begin{itemize}
    \item \textbf{[Strength 1]:} [Quantitative justification, e.g., "achieves X\% improvement over baseline Y".]
    \item \textbf{[Strength 2]:} [...]
    \item \textbf{[Strength 3]:} [Interpretability, scalability, or generalizability.]
\end{itemize}

\subsection{Limitations}
\begin{itemize}
    \item \textbf{[Limitation 1]:} [Honest statement. e.g., "Assumption 3 may not hold when..."]
    \item \textbf{[Limitation 2]:} [Computational complexity or data requirement.]
\end{itemize}

\subsection{Future Extensions}

[2–3 sentences on how the model could be extended with additional data or relaxed assumptions.]

%% ══════════════════════════════════════════════════════════════════
%%  7. CONCLUSIONS
%% ══════════════════════════════════════════════════════════════════
\section{Conclusions}

This paper presented [N] models to address the [contest problem] challenge:

\begin{enumerate}
    \item \textbf{Task 1}: [1–2 sentence result summary with specific numbers.]
    \item \textbf{Task 2}: [1–2 sentence result summary.]
    \item \textbf{Task 3}: [1–2 sentence result summary.]
\end{enumerate}

[Broader implication paragraph: what does this mean for [real-world application]?]

%% ══════════════════════════════════════════════════════════════════
%%  REFERENCES
%% ══════════════════════════════════════════════════════════════════
\bibliographystyle{apalike}
\begin{thebibliography}{99}

\bibitem{ref1}
Author, A., \& Author, B. (Year). Title of paper. \textit{Journal Name}, \textit{Vol}(Issue), pages. \url{https://doi.org/...}

\bibitem{ref2}
Author, C. (Year). \textit{Book Title}. Publisher.

% Add ≥8 references

\end{thebibliography}

%% ══════════════════════════════════════════════════════════════════
%%  APPENDICES
%% ══════════════════════════════════════════════════════════════════
\begin{appendices}

\section{Source Code}

\subsection{Data Exploration (\texttt{01\_data\_eda.py})}

\begin{lstlisting}[style=pythonstyle, caption={Data exploration and preprocessing.}]
# Insert src/01_data_eda.py here
\end{lstlisting}

\subsection{Task 1 Model (\texttt{02\_model\_problem1.py})}

\begin{lstlisting}[style=pythonstyle, caption={Task 1 model implementation.}]
# Insert src/02_model_problem1.py here
\end{lstlisting}

\end{appendices}

%% ══════════════════════════════════════════════════════════════════
%%  REPORT ON USE OF AI TOOLS（仅在实际用到 AI 工具时添加，不计入 25 页）
%% ══════════════════════════════════════════════════════════════════
%% 官方格式（COMAP Contest_AI_Policy）：每条 = 工具名 (版本/日期) + 一句话用途
%% 说明，或 Query/Output 逐字记录（篇幅较长时可用后者的摘要版）。
%% 用 scripts/ai_usage_doc.py 生成条目文本，本模板不预先填入任何内容——
%% 是否使用 AI、用了哪些、要不要加这一节，由使用者自己决定（AutoMCM_SOP.md §17）。
% \section*{Report on Use of AI Tools}
% \begin{enumerate}
% \item \textit{Tool Name} (version/date) \\
%       Query: <exact wording input to the tool> \\
%       Output: <complete output from the tool>
% \end{enumerate}

\end{document}
